\section{Codifica JPEG}
\subsection{Processo di codifica JPEG}
\begin{figure}[htbp]
	\centering
	\begin{tikzpicture}[
		% Stile per i blocchi di elaborazione
		block/.style={
			rectangle,
			draw=blue!70!black,
			fill=blue!5,
			thick,
			align=center,
			rounded corners=3pt,
			font=\footnotesize
		},
		% Stile per le frecce
		arrow/.style={
			-{Stealth[scale=1.2]},
			thick
		},
		% Stile per i segnali testuali
		signal/.style={
			align=center,
			font=\footnotesize
		},
		% Distanza standard tra i nodi
		node distance=0.7cm
	]
		% Catena di blocchi
		\node (in) at (0,0) [signal] {\textbf{Immagine}\\\textbf{YCbCr}};
		\node [block, right=of in] (cent) {Centratura\\campioni};
		\node [block, right=of cent] (blk) {Partizione\\blocchi \(8 \times 8\)};
		\node [block, right=of blk] (dct) {DCT 2D};
		\node [block, right=of dct] (quant) {Quantizzazione};
		\node [block, right=of quant] (comp) {Codifica\\Lossless};
		\node [signal, right=of comp] (out) {\textbf{Flusso}\\\textbf{JPEG}};

		% Collegamenti ed etichette
		\draw [arrow] (in) -- node[signal, above=0.1cm] {} (cent);
		\draw [arrow] (cent) -- node[signal, above=0.1cm] {} (blk);
		\draw [arrow] (blk) -- node[signal, above=0.1cm] {} (dct);
		\draw [arrow] (dct) -- node[signal, above=0.1cm] {} (quant);
		\draw [arrow] (quant) -- node[signal, above=0.1cm] {} (comp);
		\draw [arrow] (comp) -- node[signal, above=0.1cm] {} (out);
	\end{tikzpicture}
\end{figure}

\noindent
La codifica JPEG processa immagini in formato YCbCr (o in B/N), per cui il primo passo di codifica rimane la conversione
da RGB a YCbCr con relativo sottocampionamento come visto in precedenza.

La codifica JPEG lavora solo su un canale, per cui in caso di immagini multicanale (YCbCr) ogni canale viene codificato
separatamente. Il processo di codifica per un canale è composto da cinque fasi:
\begin{enumerate}
	\item centratura dei valori dei campioni dei pixel di un canale: si sottrae il valore 128 a tutti i campioni dei
	pixel di un canale, in modo da portare i valori dall'intervallo \([0, 255]\) all'intervallo \([-128, 127]\) per
	far lavorare meglio la DCT
	\item partizione dell'immagine in blocchi di dimensione \(8 \times 8\) pixel: la dimensione è stata scelta come
	compromesso tra complessità computazionale (blocchi grandi richiedono più risorse) e la qualità della compressione
	(blocchi piccoli non catturano strutture globali) in base ai calcolatori e alle risoluzioni diffusi quanto è stato
	definito lo standard (1991)
	\item calcolo dei coefficienti della matrice DCT per ogni blocco (approfondito in seguito)
	\item quantizzazione dei coefficienti DCT (approfondito in seguito)
	\item codifica lossless dei coefficienti quantizzati (approfondito in seguito)
\end{enumerate}

\noindent
Lo standard definisce la sintassi del formato, ma non pone vincoli sugli algoritmi implementativi, permettendo
lo stesso la concorrenza tra diversi produttori di codec.

\subsection{Discrete Cosine Transform - DCT}
\subsubsection*{Trasformate sparsificanti per generare residui a bassa entropia}
L'obiettivo di una codifica lossy è quello di eseguire una nuova quantizzazione dei dati per comprimere l'informazione
memorizzata. Il problema è che se si esegue una riquantizzazione diretta sui campioni, la qualità del segnale digitale
degrada molto velocemente in quanto i campioni non sono sparsi e hanno tutti approssimativamente la stessa importanza.
È possibile, però, applicare trasformate lineari sparsificanti che rappresentano il segnale in un nuovo dominio in cui
i campioni sono più sparsi e su cui è possibile eseguire una quantizzazione più aggressiva per eliminare i campioni meno
importanti senza degradare troppo la qualità del segnale. La trasformata deve essere anche facilmente invertibile per
poter tornare al dominio dei campioni originali.

\subsubsection*{Trasformate ortogonali preservano le metriche di prestazione}
Tra tutte le trasformate sparsificanti, si preferiscono quelle ortogonali per cui detta \(A\) la matrice di trasformazione
si ha che \(A^{-1} = A^T\). La proprietà di queste trasformazioni è che preservano la distanza euclidea tra due segnali,
e di conseguenza anche le metriche di prestazione come il MSE e il PSNR. Per cui una strategia di compressione ottimale
(che massimizza il PSNR nella nuova rappresentazione) è anche ottimale nel dominio dei campioni originali.

Di seguito si illustra il processo di trasformazione (1), quantizzazione (2) e ricostruzione (3) del segnale e si dimostra
che l'MSE si conserva tra il dominio dei campioni e quello dei coefficienti trasformati
\[\begin{array}{c c c}
	0. & \vec{x} = [s(n\!+\!1), s(n\!+\!2) \ldots s(n\!+\!N)]^T \in \mathbb{R}^N
	& \quad \text{segnale originale in notazione vettoriale}
	\\[5pt]
	1. & \vec{y} = A \; \vec{x}
	& \quad \text{applicazione della trasformata ortogonale}
	\\[5pt]
	2. & \tilde{y} = [Q(y(1), \Delta_1), Q(y(2), \Delta_2), \ldots, Q(y(N), \Delta_N)]
	& \quad \text{quantizzazione dei coefficienti trasformati}
	\\[5pt]
	3. & \tilde{x} = A^{-1} \; \tilde{y}
	& \quad \text{ricostruzione del segnale originale}
\end{array}\]
\vspace{-10pt}
\begin{align*}
	\text{MSE}(\vec{x}, \tilde{x}) &= \frac{1}{N} \; || \, \vec{x} - \tilde{x} \, ||^2 = \frac{1}{N} \; (\vec{x} - \tilde{x})^T \cdot (\vec{x} - \tilde{x}) = \frac{1}{N} \; (A^T \vec{y} - A^T \tilde{y})^T \cdot (A^T \vec{y} - A^T \tilde{y}) = \\
	&= \frac{1}{N} \; (\vec{y} - \tilde{y})^T \cdot A \cdot A^T \cdot (\vec{y} - \tilde{y}) = \frac{1}{N} \; || \, \vec{y} - \tilde{y} \, ||^2 = \text{MSE}(\vec{y}, \tilde{y})
\end{align*}

\subsubsection*{Trasformate bidimensionali applicate alle immagini}
Le trasformate ortogonali sono state pensate per lavorare su vettori, per cui per applicarle alle immagini (matrici) è
necessario prima applicarle alle righe e poi alle colonne della matrice dei campioni come segue. Si nota applicando la
trasformata prima sulle righe e poi sulle colonne o viceversa, il risultato non cambia.
\[Y = A \, X \, A^T = \underbrace{\, (A \, X) \,}_{\substack{\text{trasformata} \\ \text{delle colonne}}} \!\!\!\! A^T = A \underbrace{\, (A \,X^T)^T \,}_{\substack{\text{trasformata} \\ \text{sulle righe}}}\]
\begin{itemize}
	\item \(Z = A \, X\) trasformata delle colonne, il risultato è una matrice le cui colonne corrispondono al prodotto tra
	la matrice di trasformazione \(A\) e le colonne di \(X\)
	\item \(Y = Z \, A^T = (A \, Z^T)^T\) trasformata sulle righe, il risultato è una matrice le cui righe corrispondono
	al prodotto tra la matrice di trasformazione \(A\) e le righe di \(Z\)
\end{itemize}

\subsubsection*{Trasformata di Fourier e Discrete Cosine Transform applicata alle immagini}
Si nota che la trasformata di Fourier è una buona candidata per essere una trasformata sparsificante e ortogonale. In pratica
rappresenta l'immagine come combinazione di sinusoidi orizzontali e verticali. Le basse frequenze corrispondono a variazioni
lente del colore dei pixel (ad esempio superfici uniformi), mentre le alte frequenze corrispondono a variazioni rapide (ad
esempio zone ricche di dettagli). Siccome le immagini in generale sono composte principalmente da basse frequenze, la
trasformata di Fourier riesce a sparsificare bene i dati.

Nella compressione JPEG, però, si usa la Discrete Cosine Transform (DCT) che è una variante discreta della trasformata di Fourier
basata solo sulle funzioni coseno, restituisce solo coefficienti reali positivi e permette una migliore sparsificazione.

\subsubsection*{Interpretazione della DCT bidimensionale a blocchi}
Nell'algoritmo JPEG, si applica la DCT su blocchi di immagine di dimensione \(8 \times 8\) pixel, ottenendo una matrice di
coefficienti DCT di dimensione \(8 \times 8\) per ogni blocco. La scelta della dimensione del blocco è un compromesso tra
complessità e qualità della compressione.

Ogni coefficiente della matrice DCT rappresenta la correlazione (o prodotto interno) tra il blocco di pixel dell'immagine
e un particolare blocco della DCT costruito ad hoc per rappresentare un particolare pattern di frequenza. Essendo i blocchi
della DCT ordinati in frequenza crescente, le matrici di coefficienti avranno una struttura sparsificata concentrando
i coefficienti associati alle basse frequenze in alto a sinistra e quelli associati alle alte frequenze in basso a destra.

\subsubsection*{Corrispondenza tra coefficienti DCT e le strutture dell'immagine}
Analizzando la matrice di coefficienti DCT di un'immagine si notano le seguenti strutture:
\begin{itemize}
	\item superfici uniformi fanno crescere i coefficienti a bassa frequenza, concentrati in alto a sinistra
	\item i dettagli fanno crescere i coefficienti a medio/alte frequenze, concentrati al centro
	\item bordi e contorni fanno crescere i coefficienti lungo rette passanti per l'origine in alto a sinistra
	\item strutture periodiche dell'immagine formano punti isolati di medio/alta frequenza, corrispondenti proprio alla
	frequenza della struttura periodica
\end{itemize}

\subsection{Quantizzazione dei coefficienti DCT}
\subsubsection*{Quantizzazione dei coefficienti DCT}
La quantizzazione dei coefficienti DCT avviene tramite un quantizzatore uniforme e converte i coefficienti reali della DCT
in valori interi detti indici di quantizzazione. Tali indici esprimono il fattore moltiplicativo dell'intervallo di
quantizzazione \(\Delta\) necessario per ottenere il coefficiente DCT quantizzato. In questo modo si riesce a ridurre
i valori da memorizzare e rendere la codifica lossless successiva più efficiente.
\[x_\text{ind} = Q(x, \Delta) = \text{round} \left( \frac{x}{\Delta} \right) \qquad \tilde{x} = \Delta \cdot x_\text{ind} = \Delta \cdot \text{round} \left( \frac{x}{\Delta} \right)\]

\subsubsection*{Tabella di quantizzazione}
Siccome i coefficienti della DCT hanno importanza diversa in base alla loro frequenza, è possibile utilizzare quantizzatori
diversi specifici per ogni coefficiente della matrice DCT. Si definisce, quindi, una tabella di quantizzazione \(q\) di
dimensione \(8 \times 8\) che contiene i valori di quantizzazione \(\Delta(i,j)\) specifici per i corrispettivi coefficienti
\(Y(i,j)\) della matrice DCT.

Lo standard non fornisce una tabella specifica, ma ne suggerisce una di riferimento per la luminanza e una per la crominanza
ottenute tramite test soggettivi sulla qualità percepita. Tali matrici hanno valori più bassi associati alle basse frequenze
in alto a sinistra e valori più alti associati alle alte frequenze in basso a destra.

\subsubsection*{Fattore di qualità}
Per rendere la codifica più flessibile all'utente, si definisce un fattore di qualità \(Q\) compreso tra 1 e 100 che definisce
il fattore di scala \(SF\) per la tabella di quantizzazione base \(q^*\). In questo modo l'utente può scegliere se utilizzare
compressioni più aggressive (\(Q\) basso) o più conservative (\(Q\) alto).

Siccome la tabella deve essere costituita da interi nel range \([1,256]\) per essere rappresentata in 64 byte (1 byte per
valore), a seguito dello scalamento per \(SF\) viene eseguita l'operazione di troncamento della parte decimale e se il
valore è maggiore di 256, viene sostituito con 256.
\[SF \; = \; \begin{cases}
	5000 / Q & \text{se } 0 < Q \leq 50 \\
	200 - 2Q & \text{se } 50 < Q \leq 99 \\
	1 & \text{se } Q = 100
\end{cases} \qquad\qquad\qquad q \; = \; \frac{SF}{100} \; q^*\]

\subsubsection*{Effetti della quantizzazione delle frequenze}
Per immagini con bassa qualità, spesso si verificano fenomeni di sdoppiamento dei bordi e dei contorni dell'immagine.
Questo è dovuto al fatto che comprimendo pesantemente le alte frequenze, queste non riescono più a delimitare correttamente
i bordi e le zone di alta frequenza e le oscillazioni lente si propagano anche nei pixel adiacenti. Si parla di propagazione
delle strutture a bassa frequenza.

Per fattori di compressione molto bassi, si possono notare le strutture a blocchi di \(8 \times 8\) pixel.

\subsection{Codifica lossless dei coefficienti DCT}
\subsubsection*{Zig-zag scan}
Per convertire la matrice di coefficienti DCT in un vettore, si utilizza lo zig-zag scan che parte dal coefficiente
in alto a sinistra e si prosegue per diagonali ordinando i coefficienti in frequenza crescente. Al posto di riportare
tutta la sequenza di coefficienti nulli alla fine del vettore, si posiziona un simbolo speciale EOB (End Of Block) dopo
l'ultimo coefficiente non nullo.

\subsubsection*{Codifica del coefficiente DC}
Il primo coefficiente del vettore ottenuto dallo zig-zag scan è detto coefficiente DC e rappresenta la componente continua
di frequenza nulla. Per codificarlo si utilizza una codifica predittiva usando il valore DC del blocco precedente e
il residuo viene codificato con codifica categoria/valore. La tabella di codifica dei residui DC non è definita dallo
standard per cui va indicata nel file JPEG.
\[e_{\, \text{residuo}} = x_{\, \text{DC del blocco corrente}} - p_{\, \text{DC del blocco precedente}}\]

\subsubsection*{Codifica dei blocchi AC}
I restanti coefficienti del vettore sono detti AC e rappresentano le componenti di frequenza non nulla. Per codificarli
si utilizza la codifica run-length che consiste nel creare una sequenza di simboli \((R,C)\) in cui \(R\) è il numero di
coefficienti nulli consecutivi e \(C\) è il primo coefficiente AC non nullo. Se si incontrano più di 15 coefficienti
nulli, si inserisce un simbolo speciale \((15,0)\) e il successivo simbolo calcola \(R\) a partire dal successivo
coefficiente. La sequenza termina con il simbolo \((0,0)\). La codifica dei simboli \((R,C)\) avviene in due fasi:
prima si codifica \(C\) con codifica categoria/valore \((k, v)\), poi si codifica la coppia \((R,k)\) con una codifica
a prefisso. La tabella di codifica dei blocchi AC non è definita dallo standard per cui va indicata nel file JPEG.

\newpage

\subsection{Sintassi del formato JPEG baseline (1992)}
\subsubsection*{Struttura del file JPEG}
Un file JPEG è costituito da una struttura gerarchica di segmenti, come segue:
\begin{itemize}
	\item start of image: sezione di apertura del file
	\item frame header: informazioni sulla dimensione dell'immagine, sullo spazio colore, sul campionamento dei canali
	\item frame: sezione suddivisa a sua volta in uno o più scan, uno per ogni canale di cui è composta l'immagine,
	ognuno dei quali contiene:
	\begin{itemize}[topsep=0pt]
		\item scan header: informazioni sul canale e sulle tabelle di quantizzazione
		\item scan: sezione costituita da una sequenza di segmenti che possono essere codificati e decodificati in
		parallelo, ognuno dei quali contiene:
		\begin{itemize}[topsep=0pt]
			\item segment header: informazioni sulle tabelle di codifica
			\item segment: dati effettivi codificati, suddivisi in blocchi della DCT
		\end{itemize}
	\end{itemize}
	\item end of image: sezione di chiusura del file
\end{itemize}

\subsection{Evoluzione del formati di compressione delle immagini}
\begin{itemize}[itemsep=2pt]
	\item \textbf{1992 - JPEG baseline}: \\
	primo standard JPEG (\(8 \times 8\), DCT, zig-zag scan, Huffman), ampiamente usato e supportato da tutti i browser
	e sistemi operativi per la sua semplicità e compatibilità
	\item \textbf{1996 - PNG}: \\
	formato lossless con predizione spaziale e codifica con dizionario con Huffman, usato per immagini nel web e per
	immagini con trasparenza, poco diffuso in smartphone e fotocamere
	\item \textbf{2000 - JPEG 2000}: \\
	viene sostituita la DCT con la DWT (Discrete Wavelet Transform) più adatta a rappresentare le immagini, permette una
	codifica progressiva lossy-to-lossless, usato nel cinema digitale, imaging medico e immagini satellitare
	\item \textbf{2010 - WebP}: \\
	proveniente dai codec video, basato su predizione spaziale, DCT a dimensione variabile e deblocking filter, usato per
	immagini nel web
	\item \textbf{2010 - WebP lossless}: \\
	proveniente dai codec video, versione avanzata della PNG
	\item \textbf{2015 - HEIC}: \\
	proveniente dai codec video, versione avanzata di WebP, usata nell'ecosistema Apple, molto efficiente, ma soggetta
	a royalty su tecnologie brevettate (solo aziende grosse come Apple si possono permettere di gestire la burocrazia
	per ottenere le licenze)
	\item \textbf{2019 - AVIF}: \\
	versione di HEIC e WebP senza royalty, molto diffusa in ambito web e smartphone come alternativa efficiente a JPEG
	\item \textbf{2019 - JPEG XL}: \\
	versione avanzata di JPEG 2000, con codifica progressiva lossy-to-lossless, aggiunge codifica entropica dell'intero
	file, supporto per immagini HDR e migliore gestione dei bordi e texture ad alte frequenze, dovrebbe diventare il nuovo
	standard JPEG, ma il suo supporto è ancora limitato
	\item \textbf{2025 - JPEG AI}: \\
	sfrutta le codifiche basate su reti neurali (DNN e autoencoder) che utilizzano approcci di deep learning per estrapolare
	le caratteristiche principali dell'immagine e codificarle in maniera efficiente in uno spazio latente, riescono a
	modellare meglio strutture non lineari riducendo parecchio il bitrate, esistono due varianti: una ad alta complessità
	(richiede GPU e NPU) ed una a bassa complessità (sufficiente una CPU, per smartphone), potrebbe diventare il futuro
	standard JPEG
\end{itemize}
